\chapter{Аудит общего сродства и температурного якоря}
\label{ch:v10-two-reservoir-common-affinity-audit}

\section{Что осталось после происхождения разности}

Предыдущий гейт вывел
\begin{equation}
 a_c-a_h=\log2,
\end{equation}
но сохранил орбиту $(a_h,a_c)\mapsto(a_h+s,a_c+s)$. Проверим десять
кандидатов на происхождение общего сдвига и физической температуры.
Контракт требует правильного безразмерного типа, внутренней доступности,
типизации в обе ванны, выбора общего сдвига, совместимости с положительным
CPTP-каналом и независимой единицы энергии--температуры.

\section{Ближайший кандидат: пути длины один и два}

Ориентированный хопфовский цикл имеет ребро силы $F_e=\log2$. Поэтому пути
длины один и два дают точную пару
\begin{equation}
 \binom{a_h}{a_c}
 =\binom{1}{2}F_e
 =\binom{\log2}{2\log2}.
 \label{eq:v10-common-affinity-hopf-path-candidate}
\end{equation}
Это воспроизводит обе использованные KMS-величины без нового числа.
Однако текущий родитель не содержит морфизма, назначающего горячей ванне
однорёберный путь, а холодной --- двухрёберный. Без такого назначения
формула остаётся совместимым чтением уже выбранной пары.

Ту же близость даёт цепной оператор степени с уровнями $1$ и $2$, но его
целочисленный разрыв не задаёт основание $\log2$ и также не имеет
типизированного подъёма к двум текущим резервуарам.

\section{Остальные кандидаты}

Нормировка Гиббса, равенство нисходящих вероятностей, стационарная мера и
энтропийное производство вычисляют следствия заданных скоростей, но не
выбирают их общий сдвиг. Энергия часов и спектральный cutoff сохраняют
абсолютную масштабную орбиту. Размерностная трансмутация остаётся условной,
а наблюдаемая температура является внешним целевым входом.

Матрица десяти кандидатов по шести критериям имеет полный ранг $6$, но
\begin{equation}
 N_{\rm pass}=0/10,\qquad N_{\rm max}=4/6.
 \label{eq:v10-common-affinity-candidate-ledger}
\end{equation}

\section{Безразмерное сродство не является температурой}

Даже если пара $(a_h,a_c)$ задана, она фиксирует только произведения
$a_r=\beta_r\Delta$. Для общего энергетического разрыва линейная карта на
$(\log\beta_h,\log\beta_c,\log\Delta)$ имеет
\begin{equation}
 \rank C_T=2,\qquad \operatorname{nullity}C_T=1,\qquad
 C_T(-1,-1,1)^T=0.
 \label{eq:v10-common-affinity-temperature-kernel}
\end{equation}
Следовательно, обе обратные температуры контрмасштабируются с общей
энергетической щелью. Независимый разрыв повысил бы ранг до $3$, но текущий
корпус его не выбирает.

\begin{theorem}[Точный кандидат хопфовских длин]
Пути длины один и два в равнорёберном хопфовском цикле точно воспроизводят
безразмерную пару сродств $(\log2,2\log2)$ и сохраняют ранее выведенную
разность $\log2$.
\end{theorem}

\begin{nogocriterion}[Совпадение путевых длин не является назначением ванн]
Без морфизма, который выводит соответствие
«горячая ванна $\leftrightarrow$ один шаг» и
«холодная ванна $\leftrightarrow$ два шага», путевое чтение не считается
происхождением общего сдвига. Даже после такого морфизма абсолютные
температуры требуют независимой энергетической щели.
\end{nogocriterion}

\begin{protocol}\begin{enumerate}
\item проверено кандидатов: \textbf{$10$};
\item покрытие критериев: \textbf{ранг $6/6$};
\item максимальная оценка: \textbf{$4/6$};
\item полных проходов: \textbf{$0/10$};
\item общий сдвиг сродств: \textbf{$0/1$};
\item физическая температура: \textbf{$0/1$};
\item следующий гейт: \textbf{родитель назначения хопфовских путей ваннам}.
\end{enumerate}\end{protocol}

Машинный сертификат сохранён в
\path{s2t_v10_cell_birth_four_volume_induced_newton_breathing_anomaly_two_reservoir_common_affinity_temperature_anchor_candidate_audit_gate_results.json}.